\documentclass[12pt, letterpaper]{article}

% ── Paquetes ──────────────────────────────────────────────────────────────
\usepackage[utf8]{inputenc}
\usepackage[T1]{fontenc}
\usepackage[spanish]{babel}
\usepackage{geometry}
\usepackage{amsmath, amssymb, amsthm}
\usepackage{booktabs}
\usepackage{array}
\usepackage{microtype}
\usepackage{setspace}
\usepackage{parskip}
\usepackage{titlesec}
\usepackage{fancyhdr}
\usepackage{enumitem}
\usepackage{caption}
\usepackage{footnote}
\usepackage{hyperref}

% ── Tipografía estilo NYT (serif clásico) ─────────────────────────────────
\usepackage{lmodern}

% ── Geometría ─────────────────────────────────────────────────────────────
\geometry{
  top    = 2.8cm,
  bottom = 3.0cm,
  left   = 3.2cm,
  right  = 3.2cm,
}

% ── Interlineado tipo periódico ───────────────────────────────────────────
\setstretch{1.18}

% ── Encabezado y pie ─────────────────────────────────────────────────────
\pagestyle{fancy}
\fancyhf{}
\fancyhead[C]{\small\textsc{Econometría · Gujarati — Ejercicio 10.3}}
\fancyfoot[C]{\small\thepage}
\renewcommand{\headrulewidth}{0.4pt}
\renewcommand{\footrulewidth}{0pt}

% ── Títulos de sección estilo NYT ─────────────────────────────────────────
\titleformat{\section}
  {\large\bfseries\scshape}{\thesection.}{0.6em}{}[\vspace{-4pt}\rule{\linewidth}{0.4pt}]

\titleformat{\subsection}
  {\normalsize\bfseries}{\thesubsection.}{0.5em}{}

\titleformat{\subsubsection}
  {\normalsize\itshape}{}{0em}{}

% ── Teoremas ─────────────────────────────────────────────────────────────
\newtheorem{obs}{Observación}
\newtheorem{prop}{Proposición}

% ── Comandos de conveniencia ──────────────────────────────────────────────
\newcommand{\betahat}{\hat{\beta}}
\newcommand{\Rss}{SCR}
\newcommand{\Rtwo}{R^{2}}

% ─────────────────────────────────────────────────────────────────────────
\begin{document}

% ── Portada ───────────────────────────────────────────────────────────────
\begin{center}
  {\LARGE\bfseries Mortalidad Infantil, PIB, Alfabetización}\\[6pt]
  {\LARGE\bfseries y Tasa de Fecundidad Total}\\[10pt]
  {\large\itshape Un análisis de multicolinealidad y especificación}\\[14pt]
  \rule{0.55\linewidth}{0.6pt}\\[8pt]
  {\normalsize Ejercicio~10.3 — \textit{Econometría}, Gujarati \& Porter, 5.ª ed.}\\[4pt]
  {\small\textsc{Análisis detallado con desarrollo matricial y apéndice técnico}}
\end{center}

\vspace{10pt}

% ── Resumen ejecutivo ─────────────────────────────────────────────────────
\noindent\textbf{Resumen.}\quad
El ejemplo~8.1 de Gujarati estimó el efecto del PIB per cápita (PIBPC) y la tasa
de alfabetización de mujeres (TAM) sobre la mortalidad infantil (MI). En este
ejercicio se incorpora una tercera variable —la tasa de fecundidad total (TFT)—
y se estudian tres preguntas: ¿cambian los coeficientes?, ¿agrega valor la TFT?,
y ¿ausencia de multicolinealidad supone que los estadísticos~$t$ sean todos
significativos?

\vspace{6pt}

% ══════════════════════════════════════════════════════════════════════════
\section{Planteamiento del modelo}

\subsection{Modelo base (ecuación 8.1.4)}

\begin{equation}
  \widehat{\text{MI}}_i
    = \betahat_1
    + \betahat_2\,\text{PIBPC}_i
    + \betahat_3\,\text{TAM}_i
    + u_i
  \label{eq:base}
\end{equation}

\noindent donde la variable dependiente es la tasa de mortalidad infantil por cada
mil nacidos vivos, PIBPC es el producto interno bruto per cápita en dólares, y TAM
es el porcentaje de mujeres adultas con algún grado de alfabetización.

\subsection{Modelo ampliado con TFT}

Al agregar la tasa de fecundidad total (número promedio de hijos por mujer) el
modelo pasa a ser:

\begin{equation}
  \widehat{\text{MI}}_i
    = \betahat_1
    + \betahat_2\,\text{PIBPC}_i
    + \betahat_3\,\text{TAM}_i
    + \betahat_4\,\text{TFT}_i
    + u_i
  \label{eq:ampliado}
\end{equation}

% ══════════════════════════════════════════════════════════════════════════
\section{Resultados de las regresiones}

\subsection{Tabla comparativa de estimaciones}

\begin{table}[ht]
\centering
\caption{Resultados de MCO — variable dependiente: MI}
\label{tab:resultados}
\small
\begin{tabular}{lrrrr}
\toprule
\textbf{Variable} & \textbf{Coeficiente} & \textbf{Error estándar}
                  & \textbf{Estadístico $t$} & \textbf{$p$-valor} \\
\midrule
\multicolumn{5}{l}{\textit{Modelo base (ec.~\ref{eq:base}) — sin TFT}} \\[2pt]
Constante   & ---        & ---       & ---        & ---    \\
PIBPC       & $-0.0056$  & ---       & ---        & ---    \\
TAM         & $-2.2316$  & ---       & ---        & ---    \\
\midrule
\multicolumn{5}{l}{\textit{Modelo ampliado (ec.~\ref{eq:ampliado}) — con TFT}} \\[2pt]
Constante   & $168.3067$ & $32.8917$ & $5.1170$  & $0.0000$ \\
PIBPC       & $-0.0055$  & $0.0019$  & $-2.9343$ & $0.0477$ \\
TAM         & $-1.7680$  & $0.2480$  & $-7.1287$ & $0.0000$ \\
TFT         & $12.8686$  & $4.1905$  & $3.0709$  & $0.0032$ \\
\bottomrule
\end{tabular}

\vspace{4pt}
\begin{tabular}{lr@{\quad}lr}
\toprule
\multicolumn{4}{l}{\textit{Estadísticas de ajuste — modelo ampliado}} \\
\midrule
$R^2$                     & $0.7474$ & Media de MI    & $141.50$  \\
$\bar{R}^2$               & $0.7347$ & Desv. est. MI  & $75.9781$ \\
Error estándar de regr.   & $39.1313$& Criterio AIC   & $10.2322$ \\
SCR (suma cuad. resid.)   & $91\,875.38$ & Criterio SC & $10.3671$ \\
Log-verosimilitud         & $-323.43$& $F$-estadístico& $59.17$   \\
Durbin-Watson             & $2.1703$ & $p(F)$         & $0.0000$  \\
\bottomrule
\end{tabular}
\end{table}

% ══════════════════════════════════════════════════════════════════════════
\section{Respuesta a los incisos}

% ── a) ────────────────────────────────────────────────────────────────────
\subsection{Inciso (a): comparación de coeficientes y explicación de cambios}

Al comparar la ecuación~(\ref{eq:base}) con la~(\ref{eq:ampliado}) se observan
dos cambios principales.

\medskip
\noindent\textbf{Cambio 1 — Coeficiente de TAM.}\quad
En el modelo base $\betahat_3 \approx -2.2316$, mientras que en el modelo ampliado
$\betahat_3 = -1.7680$. El coeficiente se reduce en valor absoluto. Esto ocurre
porque \emph{parte del efecto que antes se atribuía a TAM ahora es captado por TFT}:
la alfabetización y la fecundidad están correlacionadas —países con mayor
alfabetización femenina tienden a tener menor fecundidad—, por lo que la omisión
de TFT en el modelo base producía un sesgo al alza (en valor absoluto) en el
estimador de TAM. Este fenómeno se denomina \textbf{sesgo por variable omitida}
(\textit{omitted variable bias}):

\begin{equation}
  \text{plim}\,\hat{\beta}_3^{\text{base}}
  = \beta_3
  + \beta_4 \cdot \frac{\text{Cov}(\text{TFT},\text{TAM})}{\text{Var}(\text{TAM})}
  \label{eq:sesgo}
\end{equation}

\noindent Si $\beta_4 > 0$ (TFT aumenta MI) y $\text{Cov}(\text{TFT},\text{TAM}) < 0$
(mayor alfabetización implica menor fecundidad), el término de sesgo es negativo,
lo que magnifica $|\hat{\beta}_3^{\text{base}}|$ respecto al verdadero $\beta_3$.

\medskip
\noindent\textbf{Cambio 2 — Coeficiente de PIBPC.}\quad
El coeficiente de PIBPC pasa de $-0.0056$ a $-0.0055$, prácticamente sin cambios.
Esto sugiere que PIBPC y TFT no están tan correlacionadas, o que la correlación
parcial de TFT con PIBPC —controlando por TAM— es pequeña.

\medskip
\noindent\textbf{Cambio 3 — $R^2$.}\quad
El $R^2$ sube de forma modesta al agregar TFT. Es de esperar: agregar un regresor
nunca puede reducir el $R^2$ en muestra. La pregunta relevante es si el ajuste
\emph{fuera de muestra} mejora, lo cual conduce al inciso (b).

% ── b) ────────────────────────────────────────────────────────────────────
\subsection{Inciso (b): ¿vale la pena agregar TFT?}

Para responder formalmente usamos el \textbf{test $F$ parcial} (o test de
Wald), que contrasta si el coeficiente de la variable añadida es cero en
la población:

\begin{equation}
  H_0: \beta_4 = 0 \qquad \text{vs} \qquad H_1: \beta_4 \neq 0
  \label{eq:hipotesis_b}
\end{equation}

\subsubsection{Criterio 1 — Estadístico $t$}

El estadístico de prueba es:
\begin{equation}
  t = \frac{\hat{\beta}_4}{\text{ee}(\hat{\beta}_4)}
    = \frac{12.8686}{4.1905}
    = 3.0709
  \label{eq:t_TFT}
\end{equation}

\noindent Con $p$-valor $= 0.0032 < 0.05$, se rechaza $H_0$. La TFT
\textbf{es estadísticamente significativa} al 1\%.

\subsubsection{Criterio 2 — $\bar{R}^2$ (penaliza por número de regresores)}

\begin{equation}
  \bar{R}^2 = 1 - \frac{n-1}{n-k}\left(1 - R^2\right)
  \label{eq:R2adj}
\end{equation}

\noindent El $\bar{R}^2$ sube al agregar TFT (de $\approx 0.72$ a $0.7347$),
confirmando que la ganancia en ajuste supera la penalización por el parámetro
extra.

\subsubsection{Criterio 3 — Criterio de información de Akaike (AIC)}

\begin{equation}
  \text{AIC} = -\frac{2\ell}{n} + \frac{2k}{n}
  \label{eq:AIC}
\end{equation}

\noindent donde $\ell$ es el logaritmo de la verosimilitud y $k$ es el número de
parámetros. Un AIC menor indica mejor modelo ajustando por complejidad.

\medskip
\noindent\textbf{Conclusión del inciso (b).}\quad Sí vale la pena agregar TFT.
Los tres criterios son consistentes: el estadístico $t$ es significativo al 1\%,
el $\bar{R}^2$ mejora, y el AIC desciende. Además, la TFT tiene una interpretación
económica clara: a mayor número de hijos por mujer, mayor presión sobre los recursos
familiares y el sistema de salud, elevando la mortalidad infantil.

% ── c) ────────────────────────────────────────────────────────────────────
\subsection{Inciso (c): ¿significancia individual garantiza ausencia de multicolinealidad?}

\textbf{No.} Esta es quizás la lección más importante del ejercicio. Que todos los
coeficientes sean individualmente significativos \emph{no implica} ausencia de
multicolinealidad. Los dos conceptos son distintos:

\begin{itemize}[leftmargin=1.6em, itemsep=4pt]
  \item \textbf{Significancia individual} ($t$-test): mide si un coeficiente es
    distinguible de cero \emph{dado que los demás regresores están en el modelo}.
    Con muestras grandes o varianza elevada de los regresores, incluso coeficientes
    con correlación moderada pueden ser significativos.

  \item \textbf{Multicolinealidad}: describe la correlación entre regresores. Su
    efecto es \emph{inflar los errores estándar}, pero no necesariamente hasta el
    punto de hacer los coeficientes no significativos —depende del tamaño muestral
    y de la varianza total.
\end{itemize}

\subsubsection{Indicadores formales de multicolinealidad}

\noindent\textbf{(i) Factor de inflación de varianza (VIF):}

\begin{equation}
  \text{VIF}_j = \frac{1}{1 - R_j^2}
  \label{eq:VIF}
\end{equation}

\noindent donde $R_j^2$ es el $R^2$ de la regresión auxiliar de $X_j$ sobre los
demás regresores. Un VIF $> 10$ indica multicolinealidad severa.

\medskip
\noindent\textbf{(ii) Tolerancia:}
\begin{equation}
  \text{TOL}_j = 1 - R_j^2 = \frac{1}{\text{VIF}_j}
  \label{eq:TOL}
\end{equation}

\noindent Valores de TOL próximos a~0 señalan alta colinealidad.

\medskip
\noindent\textbf{(iii) Número de condición de la matriz $\mathbf{X'X}$:}
\begin{equation}
  \kappa = \sqrt{\frac{\lambda_{\max}}{\lambda_{\min}}}
  \label{eq:condicion}
\end{equation}

\noindent donde $\lambda_{\max}$ y $\lambda_{\min}$ son el mayor y menor valor
propio de $\mathbf{X'X}$ (o de su versión estandarizada). Si $\kappa > 30$, la
colinealidad es problemática.

\subsubsection{Por qué los $t$ pueden ser significativos bajo colinealidad}

Considere el caso de dos regresores $X_2$, $X_3$ con correlación $r_{23}$.
El error estándar del estimador MCO de $\beta_2$ es:

\begin{equation}
  \text{ee}(\hat{\beta}_2)
    = \frac{\hat{\sigma}}{\sqrt{\sum x_{2i}^2}}
      \cdot \frac{1}{\sqrt{1 - r_{23}^2}}
  \label{eq:ee_multicol}
\end{equation}

\noindent El factor $1/\sqrt{1-r_{23}^2} \geq 1$ infla el error estándar. Pero si
$\hat{\sigma}$ es pequeño (buen ajuste global) y $\sum x_{2i}^2$ es grande
(mucha variación en $X_2$), el $t$ puede seguir siendo grande aun con
correlación alta entre regresores. \textbf{En el presente caso, el $R^2 = 0.747$
y la muestra es moderada}, lo que permite que los tres $t$ sean significativos
a pesar de que PIBPC, TAM y TFT están correlacionadas entre sí.

\medskip
\noindent\textbf{Conclusión del inciso (c).}\quad La significancia individual es
condición \emph{suficiente} para rechazar que un coeficiente sea cero, pero
\emph{no es condición necesaria ni suficiente} para descartar multicolinealidad.
El diagnóstico correcto requiere calcular los VIF y el número de condición de la
matriz de regresores.

% ══════════════════════════════════════════════════════════════════════════
\section{Síntesis}

\begin{enumerate}[leftmargin=1.8em, itemsep=6pt]
  \item Agregar TFT \textbf{modifica} el coeficiente de TAM porque corrige el
    sesgo por variable omitida que existía en el modelo base. PIBPC permanece
    prácticamente invariante.

  \item TFT \textbf{debe incluirse}: es significativa al 1\%, mejora el $\bar{R}^2$
    y reduce el AIC, además de estar respaldada por teoría económica y demográfica.

  \item La significancia individual de todos los coeficientes \textbf{no descarta
    multicolinealidad}. El diagnóstico debe basarse en VIF, tolerancia y número
    de condición.
\end{enumerate}

% ══════════════════════════════════════════════════════════════════════════
\appendix
\section{Apéndice técnico: álgebra matricial del estimador MCO}

\subsection{Notación}

Sea $\mathbf{y}$ el vector $n\times 1$ de observaciones de MI, y
$\mathbf{X}$ la matriz $n\times k$ de regresores (incluyendo la columna de unos
para el intercepto). El modelo es:

\begin{equation}
  \mathbf{y} = \mathbf{X}\boldsymbol{\beta} + \mathbf{u},
  \qquad \mathbf{u} \sim \mathcal{N}(\mathbf{0},\, \sigma^2 \mathbf{I}_n)
  \label{eq:modelo_matricial}
\end{equation}

\subsection{Estimador MCO}

Minimizando $\mathbf{u'u} = (\mathbf{y}-\mathbf{X}\boldsymbol{\beta})'(\mathbf{y}-\mathbf{X}\boldsymbol{\beta})$
respecto a $\boldsymbol{\beta}$, las condiciones de primer orden dan las
\textbf{ecuaciones normales}:

\begin{equation}
  \mathbf{X'X}\hat{\boldsymbol{\beta}} = \mathbf{X'y}
  \label{eq:normales}
\end{equation}

\noindent Siempre que $\mathbf{X'X}$ sea invertible (no hay multicolinealidad
perfecta), la solución única es:

\begin{equation}
  \hat{\boldsymbol{\beta}} = (\mathbf{X'X})^{-1}\mathbf{X'y}
  \label{eq:betahat}
\end{equation}

\subsection{Varianza del estimador}

\begin{equation}
  \text{Var}(\hat{\boldsymbol{\beta}})
    = \sigma^2 (\mathbf{X'X})^{-1}
  \label{eq:var_beta}
\end{equation}

\noindent En la práctica $\sigma^2$ se estima con:

\begin{equation}
  \hat{\sigma}^2
    = \frac{\mathbf{e'e}}{n - k}
    = \frac{\text{SCR}}{n - k}
  \label{eq:sigma2}
\end{equation}

\noindent donde $\mathbf{e} = \mathbf{y} - \mathbf{X}\hat{\boldsymbol{\beta}}$
son los residuos. El error estándar del coeficiente $j$-ésimo es:

\begin{equation}
  \text{ee}(\hat{\beta}_j)
    = \hat{\sigma}\sqrt{c_{jj}}
  \label{eq:ee_j}
\end{equation}

\noindent con $c_{jj}$ el elemento $(j,j)$ de $(\mathbf{X'X})^{-1}$.

\subsection{Efecto de la multicolinealidad sobre $(\mathbf{X'X})^{-1}$}

Cuando dos columnas de $\mathbf{X}$ están altamente correlacionadas, la matriz
$\mathbf{X'X}$ se aproxima a ser singular: su determinante se acerca a cero y
los elementos diagonales de $(\mathbf{X'X})^{-1}$ se vuelven muy grandes. Por
la ecuación~(\ref{eq:var_beta}), esto infla las varianzas de los estimadores:

\begin{equation}
  \text{Var}(\hat{\beta}_j)
    = \sigma^2 \cdot c_{jj}
    \quad \xrightarrow[\text{colinealidad}]{\text{alta}} \quad \infty
  \label{eq:inflacion}
\end{equation}

\noindent Intuitivamente, si $X_2 \approx \lambda X_3$, los datos no contienen
información suficiente para separar el efecto de $X_2$ del de $X_3$.

\subsection{Fórmula del VIF en el caso de tres regresores}

Suponga el modelo con PIBPC ($X_2$), TAM ($X_3$) y TFT ($X_4$). El VIF de $X_2$
se obtiene regresando $X_2$ sobre $X_3$ y $X_4$:

\begin{equation}
  X_{2i} = \gamma_1 + \gamma_3 X_{3i} + \gamma_4 X_{4i} + v_i
  \label{eq:auxiliar}
\end{equation}

\begin{equation}
  \text{VIF}_2 = \frac{1}{1 - R_{2\cdot 34}^2}
  \label{eq:VIF2}
\end{equation}

\noindent donde $R_{2\cdot 34}^2$ es el $R^2$ de la regresión~(\ref{eq:auxiliar}).
Se repite análogamente para $X_3$ y $X_4$.

\subsection{Test $F$ de significancia conjunta}

Para contrastar $H_0\!: \beta_2 = \beta_3 = \beta_4 = 0$ (excluir todos los
regresores, salvo el intercepto):

\begin{equation}
  F = \frac{(R^2/k-1)}{(1-R^2)/(n-k)}
    \sim F_{k-1,\, n-k}
  \label{eq:F_global}
\end{equation}

\noindent Con los datos del ejercicio:
$F = \frac{0.7474/3}{0.2526/60} = 59.17$, $p < 0.001$.

\subsection{Test $F$ incremental para TFT}

Sean $R_r^2$ el $R^2$ del modelo restringido (sin TFT) y $R_{ur}^2$ el del
modelo no restringido (con TFT). El test incremental es:

\begin{equation}
  F^* = \frac{(R_{ur}^2 - R_r^2)/q}{(1 - R_{ur}^2)/(n - k)}
  \label{eq:F_incremental}
\end{equation}

\noindent donde $q = 1$ es el número de restricciones. Bajo $H_0$ esta
estadística sigue una $F_{1,\, n-k}$. Nótese que $F^* = t^2$ cuando $q=1$:
$3.0709^2 \approx 9.43 \approx F_{1,60}$ al $p = 0.0032$.

% ──────────────────────────────────────────────────────────────────────────
\vspace{18pt}
\noindent\rule{\linewidth}{0.4pt}

\noindent{\small\textit{Nota.} Los datos corresponden a una muestra de 64 países usada
en Gujarati \& Porter (2010). Los cálculos presentados asumen errores
esféricos; para inferencia robusta a heterocedasticidad utilice errores
estándar de White.}

\end{document}
